
\subsubsection{4.1 Datasets}

\textbf{HumanEval}: 164 Python programming problems designed to evaluate functional correctness of code generation models.

\textbf{MBPP}: 500 Python programming problems covering a broader range of programming concepts.

A total of 664 samples were used, with 500 valid samples after filtering.

\subsubsection{4.2 Implementation Details}

\begin{itemize}
\item \textbf{Trace Collection}: Python's \texttt{sys.settrace()} with SIGALRM-based 5.0s timeout
\item \textbf{Token Mapping}: GPT-2 tokenizer offset\_mapping for line-to-token alignment
\item \textbf{Gradient Verification}: PyTorch autograd with explicit gradient norm computation
\item \textbf{Seeds}: 42, 123, 456 for variance estimation
\end{itemize}

\subsubsection{4.3 Conditions}

\begin{table}[htbp]
\centering
\resizebox{\columnwidth}{!}{%
\begin{tabular}{ll}
\toprule
Condition & Description \\
\midrule
none & No masking; standard PPO baseline \\
random & Random masking at matched sparsity (~81\%) \\
trace & Trace-based FGO masking \\
\bottomrule
\end{tabular}
}
\end{table}

\subsubsection{4.4 Evaluation}

H-M3 efficiency validation uses simulation-based learning curves rather than full PPO training. This represents a proof-of-concept validation; full training comparison is deferred to future work.
